\chapter{Complexity — Random Walk}
%%% generated by the blueprint pipeline from TCSlib.Complexity.Randomized.RandomWalk
\section{Overview}

Declarations of the Lean module \texttt{TCSlib.Complexity.Randomized.RandomWalk}.

\section{Declarations}

\begin{definition}[Probability of an index event]
\lean{Randomized.idxProb}
\label{Randomized.idxProb}
\leanok
Let $D$ and $k$ be natural numbers, and let $A$ be a property of lists of elements of
$	ext{Fin}(D)$. The probability of $A$ under a uniformly chosen length-$k$ sequence from
$	ext{Fin}(D)$ is the number of such sequences for which $A$ holds, divided by $D^k$.
\end{definition}

\begin{theorem}[Probability at length zero]
\lean{Randomized.idxProb_zero}
\label{Randomized.idxProb_zero}
\leanok
\uses{Randomized.idxProb}
\difficulty{3}
Let $D$ be a natural number and let $A$ be a predicate on finite sequences of symbols
from an alphabet of size $D$. The probability assigned to $A$ on sequences of length
zero is $1$ if $A$ holds for the empty sequence, and $0$ otherwise.
\end{theorem}

\begin{theorem}[Recurrence for index probability]
\lean{Randomized.idxProb_succ}
\label{Randomized.idxProb_succ}
\leanok
\uses{Randomized.idxProb}
\difficulty{5}
Let $D$ be a positive natural number, let $k$ be a natural number, and let $A$ be a
decidable property of finite lists of elements of $\text{Fin}(D)$. The probability that
a uniformly chosen list of length $k+1$ satisfies $A$ equals $1/D$ times the sum, over
all $i\in\text{Fin}(D)$, of the probability that the list obtained by prepending $i$ to
a uniformly chosen list of length $k$ satisfies $A$.
\end{theorem}

\begin{theorem}[Probability of a certain event]
\lean{Randomized.idxProb_true}
\label{Randomized.idxProb_true}
\leanok
\uses{Randomized.idxProb}
\difficulty{3}
Let $D>0$ and $k$ be integers, and choose uniformly from the length-$k$ sequences over
an alphabet of size $D$. The probability that a predicate which is true of every
sequence holds is $1$.
\end{theorem}

\begin{definition}[Undirected regular multigraph]
\lean{Randomized.RegularMultigraph}
\statementsource{Vad12}{
  pdf-bf2ce5422088-p038-b004,
  pdf-bf2ce5422088-p108-b004
}
\statementsource{arora-barak-ch07-randomized-computation}{pdf-26d95080dbde-p021-b002}
\label{Randomized.RegularMultigraph}
\leanok
A regular multigraph of degree $D$ on a type $V$ assigns to each vertex $v$ and each of
its $D$ edge slots a neighboring vertex. It is undirected in the sense that, for every
pair of vertices $u,v$, the number of edge slots at $u$ leading to $v$ equals the number
of edge slots at $v$ leading to $u$.
\end{definition}

\begin{definition}[Edge multiplicity]
\lean{Randomized.RegularMultigraph.mult}
\statementsource{Vad12}{pdf-bf2ce5422088-p038-b004}
\statementsource{arora-barak-ch07-randomized-computation}{pdf-26d95080dbde-p021-b002}
\label{Randomized.RegularMultigraph.mult}
\leanok
For a regular multigraph $G$ and vertices $u,v$, define the multiplicity from $u$ to $v$
to be the number of edge slots at $u$ whose assigned neighbor is $v$.
\end{definition}

\begin{definition}[Adjacency in a multigraph]
\lean{Randomized.RegularMultigraph.Adj}
\statementsource{Vad12}{pdf-bf2ce5422088-p039-b002}
\statementsource{arora-barak-ch07-randomized-computation}{pdf-26d95080dbde-p021-b002}
\label{Randomized.RegularMultigraph.Adj}
\leanok
For a regular multigraph $G$ and vertices $u,v$, declare $u$ adjacent to $v$ if at least
one edge slot at $u$ leads to $v$.
\end{definition}

\begin{theorem}[Symmetry of multigraph adjacency]
\lean{Randomized.RegularMultigraph.adj_symm}
\label{Randomized.RegularMultigraph.adj_symm}
\leanok
\uses{Randomized.RegularMultigraph.Adj}
\difficulty{4}
Let $G$ be a regular multigraph. If a vertex $u$ is adjacent to a vertex $v$, then $v$
is adjacent to $u$.
\end{theorem}

\begin{definition}[Neighbor average]
\lean{Randomized.RegularMultigraph.avg}
\statementsource{arora-barak-ch07-randomized-computation}{pdf-26d95080dbde-p021-b002}
\label{Randomized.RegularMultigraph.avg}
\leanok
For a regular multigraph $G$ of degree $D$, a function $f$ from its vertices to the
rationals, and a vertex $v$, define the neighbor average at $v$ to be the sum of $f$
over the neighbors assigned to the $D$ edge slots at $v$, divided by $D$.
\end{definition}

\begin{theorem}[Sum over neighbors by multiplicity]
\lean{Randomized.RegularMultigraph.sum_nbr_eq}
\label{Randomized.RegularMultigraph.sum_nbr_eq}
\leanok
\uses{Randomized.RegularMultigraph.mult}
\difficulty{4}
Let $G$ be a finite regular multigraph of degree $D$, let $u$ be a vertex, and let $h$
assign a rational number to each vertex. The sum of $h$ over the $D$ neighbor indices of
$u$ equals the sum over vertices $v$ of the number of indices leading from $u$ to $v$,
multiplied by $h(v)$.
\end{theorem}

\begin{theorem}[Symmetry of the neighbor sum]
\lean{Randomized.RegularMultigraph.sum_swap}
\label{Randomized.RegularMultigraph.sum_swap}
\leanok
\uses{Randomized.RegularMultigraph.mult, Randomized.RegularMultigraph.sum_nbr_eq}
\difficulty{4}
Let $G$ be a finite regular multigraph of degree $D$, and let $g$ be a rational-valued
function of two vertices. Summing $g(u,v)$ over every vertex $u$ and every indexed
neighbor $v$ of $u$ gives the same result as summing $g(v,u)$ over those same ordered
pairs.
\end{theorem}

\begin{theorem}[Total incoming multiplicity]
\lean{Randomized.RegularMultigraph.sum_mult_left}
\label{Randomized.RegularMultigraph.sum_mult_left}
\leanok
\uses{Randomized.RegularMultigraph.mult, Randomized.RegularMultigraph.sum_nbr_eq}
\difficulty{4}
Let $G$ be a finite regular multigraph of degree $D$, and let $v$ be a vertex. For each
vertex $u$, count the neighbor indices at which $u$ points to $v$; the sum of these
multiplicities over all $u$ is $D$.
\end{theorem}

\begin{theorem}[Uniform Stationary Mass Identity]
\lean{Randomized.RegularMultigraph.uniform_stationary}
\label{Randomized.RegularMultigraph.uniform_stationary}
\leanok
\uses{Randomized.RegularMultigraph.mult, Randomized.RegularMultigraph.sum_mult_left}
\difficulty{3}
Let $G$ be a regular multigraph of positive degree $D$ on a finite vertex set, and let
$v$ be a vertex. The sum over all vertices $u$ of $1/|V|$ times the natural-number
quotient of the number of neighbor indices at $u$ leading to $v$ by $D$ equals $1/|V|$.
\end{theorem}

\begin{theorem}[Sum of neighbor averages]
\lean{Randomized.RegularMultigraph.sum_avg}
\label{Randomized.RegularMultigraph.sum_avg}
\leanok
\uses{Randomized.RegularMultigraph.avg, Randomized.RegularMultigraph.sum_swap}
\difficulty{4}
Let $G$ be a finite regular multigraph of positive degree $D$, and let $f$ be a
rational-valued function on its vertices. The sum, over all vertices, of the average of
$f$ over the indexed neighbors of each vertex equals the sum of $f$ over all vertices.
\end{theorem}

\begin{definition}[Path hits a target]
\lean{Randomized.RegularMultigraph.hits}
\statementsource{Vad12}{pdf-bf2ce5422088-p038-b006}
\statementsource{arora-barak-ch07-randomized-computation}{
  pdf-26d95080dbde-p016-b002,
  pdf-26d95080dbde-p021-b002
}
\label{Randomized.RegularMultigraph.hits}
\leanok
For a regular multigraph $G$, target vertex $t$, starting vertex $v$, and a list of edge
slots, define whether the corresponding walk hits $t$ as follows: it does if $v=t$, or
if, after following the first listed slot from $v$, the walk specified by the remaining
slots hits $t$. The empty list hits $t$ exactly when $v=t$.
\end{definition}

\begin{theorem}[A hit implies reachability]
\lean{Randomized.RegularMultigraph.reflTransGen_of_hits}
\label{Randomized.RegularMultigraph.reflTransGen_of_hits}
\leanok
\uses{Randomized.RegularMultigraph.Adj, Randomized.RegularMultigraph.hits}
\difficulty{4}
Let $G$ be a regular multigraph, and let $s,t$ be vertices. If a finite sequence of
neighbor indices, followed from $s$, visits $t$ (including a visit at the starting
vertex), then there is a finite adjacency path from $s$ to $t$.
\end{theorem}

\begin{definition}[Hit probability]
\lean{Randomized.RegularMultigraph.hitProb}
\statementsource{Vad12}{pdf-bf2ce5422088-p038-b006}
\statementsource{arora-barak-ch07-randomized-computation}{pdf-26d95080dbde-p016-b002}
\label{Randomized.RegularMultigraph.hitProb}
\leanok
\uses{Randomized.RegularMultigraph.hits, Randomized.idxProb}
Let $G$ be a regular multigraph of degree $D$ on a vertex type $V$. For vertices $t,s\in
V$ and an integer $k\ge 0$, the hit probability is the index-sampling probability that
the predicate $G.\mathrm{hits}(t,s,w)=\mathrm{true}$ holds.
\end{definition}

\begin{definition}[Miss probability]
\lean{Randomized.RegularMultigraph.missProb}
\statementsource{arora-barak-ch07-randomized-computation}{pdf-26d95080dbde-p016-b002}
\label{Randomized.RegularMultigraph.missProb}
\leanok
\uses{Randomized.RegularMultigraph.hitProb}
Let $G$ be a regular multigraph of degree $D$ on a vertex type $V$. For vertices $t,s\in
V$ and an integer $k\ge 0$, the miss probability is one minus the hit probability for
$t$, $k$, and $s$.
\end{definition}

\begin{theorem}[Zero-step hitting probability]
\lean{Randomized.RegularMultigraph.hitProb_zero}
\label{Randomized.RegularMultigraph.hitProb_zero}
\leanok
\uses{Randomized.RegularMultigraph.hitProb, Randomized.RegularMultigraph.hits, Randomized.idxProb_zero}
\difficulty{3}
Let $G$ be a regular multigraph of degree $D$, and let $s,t$ be vertices. The
probability that a length-zero random walk starting at $s$ hits $t$ is $1$ if $s=t$ and
$0$ otherwise.
\end{theorem}

\begin{theorem}[One-step hit probability recursion]
\lean{Randomized.RegularMultigraph.hitProb_succ}
\label{Randomized.RegularMultigraph.hitProb_succ}
\leanok
\uses{Randomized.RegularMultigraph.hitProb, Randomized.RegularMultigraph.hits, Randomized.idxProb_succ, Randomized.idxProb_true}
\difficulty{4}
Let $G$ be a regular multigraph of degree $D>0$, and let $s,t$ be vertices. For every
$k\in\mathbb{N}$, the probability that a uniformly chosen sequence of $k+1$ neighbor
indices starting at $s$ hits $t$ is $1$ if $s=t$; otherwise, it is the sum of the
length-$k$ hit probabilities from the neighbors of $s$, divided by $D$.
\end{theorem}

\begin{theorem}[Recurrence for the miss probability]
\lean{Randomized.RegularMultigraph.missProb_succ}
\label{Randomized.RegularMultigraph.missProb_succ}
\leanok
\uses{Randomized.RegularMultigraph.avg, Randomized.RegularMultigraph.hitProb_succ, Randomized.RegularMultigraph.missProb}
\difficulty{4}
Let $G$ be a finite regular multigraph of degree $D>0$, and let $t,s$ be vertices and
$k$ a nonnegative integer. The miss probability after $k+1$ steps from $s$ equals $0$ if
$s=t$; otherwise it equals the average, over the neighbor slots of $s$, of the miss
probability after $k$ steps from the resulting vertex.
\end{theorem}

\begin{theorem}[Hit probability is nonnegative]
\lean{Randomized.RegularMultigraph.hitProb_nonneg}
\label{Randomized.RegularMultigraph.hitProb_nonneg}
\leanok
\uses{Randomized.RegularMultigraph.hitProb, Randomized.idxProb}
\difficulty{1}
Let $G$ be a regular multigraph, and let $s,t$ be vertices. For every natural number
$k$, the probability that a uniformly chosen sequence of $k$ neighbor indices starting
at $s$ visits $t$ is nonnegative.
\end{theorem}

\begin{theorem}[Hit probability is at most one]
\lean{Randomized.RegularMultigraph.hitProb_le_one}
\label{Randomized.RegularMultigraph.hitProb_le_one}
\leanok
\uses{Randomized.RegularMultigraph.hitProb, Randomized.RegularMultigraph.hits, Randomized.idxProb}
\difficulty{4}
Let $G$ be a regular multigraph of degree $D>0$, and let $s,t$ be vertices. For every
natural number $k$, the probability that a uniformly chosen sequence of $k$ neighbor
indices starting at $s$ visits $t$ is at most $1$.
\end{theorem}

\begin{theorem}[Bounds on the miss probability]
\lean{Randomized.RegularMultigraph.missProb_mem_Icc}
\label{Randomized.RegularMultigraph.missProb_mem_Icc}
\leanok
\uses{Randomized.RegularMultigraph.hitProb_le_one, Randomized.RegularMultigraph.hitProb_nonneg, Randomized.RegularMultigraph.missProb}
\difficulty{3}
Let $G$ be a regular multigraph of degree $D>0$, and let $t$ and $s$ be vertices. For
every nonnegative integer $k$, the probability that a length-$k$ sequence of uniformly
chosen neighbor slots, started at $s$, fails to visit $t$ lies between $0$ and $1$,
inclusive.
\end{theorem}

\begin{definition}[Decoding index blocks]
\lean{Randomized.decodeIdx}
\label{Randomized.decodeIdx}
\leanok
Let $b,D$ be natural numbers and let $e$ be a bijection from Boolean strings of length
$b$ to $	ext{Fin}(D)$. For a natural number $k$ and a Boolean list, decode $k$
successive blocks of length $b$ by applying $e$ to each block, treating any positions
beyond the list as false; the resulting list has length $k$.
\end{definition}

\begin{theorem}[Decoding a block followed by a suffix]
\lean{Randomized.decodeIdx_succ_append}
\label{Randomized.decodeIdx_succ_append}
\leanok
\uses{Randomized.decodeIdx}
\difficulty{4}
Let $b,D$ be natural numbers, let $e$ be a bijection from the Boolean strings of length
$b$ to the indices $0,\ldots,D-1$, let $k$ be a natural number, let $u$ be a Boolean
string of length $b$, and let $l$ be any finite Boolean string. Decoding $k+1$ blocks
from the concatenation of $u$ and $l$ gives first the index $e(u)$, followed by the
result of decoding $k$ blocks from $l$.
\end{theorem}

\begin{theorem}[Bit-string simulation of indexed walks]
\lean{Randomized.RegularMultigraph.randProb_hits_decodeIdx}
\label{Randomized.RegularMultigraph.randProb_hits_decodeIdx}
\leanok
\uses{Randomized.RegularMultigraph.hitProb, Randomized.RegularMultigraph.hitProb_succ, Randomized.RegularMultigraph.hitProb_zero, Randomized.RegularMultigraph.hits, Randomized.decodeIdx, Randomized.decodeIdx_succ_append, Randomized.randProb, Randomized.randProb_append, Randomized.randProb_true}
\difficulty{6}
Let $b$ be a natural number, and suppose there is a bijection from the length-$b$
Boolean strings to the $D$ neighbor indices. For any target vertex $t$, any natural
number $k$, and any starting vertex $s$, the probability that a uniformly chosen Boolean
string of length $kb$, divided into $k$ blocks and decoded into neighbor indices, visits
$t$ equals the probability that a uniformly chosen sequence of $k$ neighbor indices
visits $t$.
\end{theorem}
